\scanpage{304}
\mspage{220}

b) Je présume qu'en dehors d'une hypothèse restrictive, du type a) ou b) ci-dessus, le foncteur \inserted{pl fid.} (4.9.3) n'est pas tjrs ess. surjectif, même dans le cas particulier où $J$ est une catégorie finie, p.ex. $J=B_G$, $G=\{\pm1\}$. Pour un contre-exemple dans le cas très voisin où on prend $M=\operatorname{Ind}(C)$ (au lieu de $\operatorname{Ind}_{\pi_0}(C)$), et où on considère le foncteur \inserted{(pl. fid. si $J$ est une catégorie équiv. à une cat. \uline{finie})}

\[
\operatorname{Ind}(\underline{\operatorname{Hom}}(J,C))\longrightarrow\underline{\operatorname{Hom}}(J,\operatorname{Ind}(C)),
\]

je renvoie à l'exercice SGA 4 I 8.8.8.

\unsure{3)|c)} L'application principale du th. 4.9.3. consiste en le

\uline{Corollaire} 4.9.5. Soient $M$ une catégorie \uline{accessible} (4.8.7), $J$ une catégorie équivalente à une petite catégorie. Alors la catégorie \inserted{$M(J^\circ)=$} $\underline{\operatorname{Hom}}(J,M)$ est accessible. Plus précisément, soient $\pi_0$ un cardinal \uline{utile} pour $M$ (i.e tel que $\operatorname{Ind}_{\pi_0}(M_{\pi_0})\xrightarrow{\approx} M$, \deleted{\unsure{) et}} \deleted{soit $\pi$} i.e. \unsure{tout} cardinal tel que $M$ soit $\pi_0$-accessible), soit $\pi$ un cardinal \inserted{infini} \deleted{\unsure{tel}} que

\[
\pi^{\pi_0}=\pi\geq\operatorname{card}\operatorname{red}J,\quad\overbrace{\deleted{\text{(et }\pi\unsure{\geq|>}\operatorname{card}\operatorname{red} M_{\pi_0}}}^{\inserted{\ast}}\quad\deleted{\text{p. ex.}}\inserted{\text{i.e.}}\quad\pi=\deleted{\unsure{C}}c^{\pi_0},\quad\text{où }\deleted{\unsure{\pi}}
\]

\margin{(On a donc $\pi\deleted{\unsure{\geq}}=\pi^{\pi_0}\geq 2^{\pi_0}>\pi_0$)}% slanted

\deleted{$\pi'_{\unsure{1|g}}=\operatorname{Sup}(\pi_0,\operatorname{card}\operatorname{red}J)$ \unsure{est tel que}} $c$ est n'importe quel cardinal $c\geq\operatorname{Sup}(\pi_0,\operatorname{card}\operatorname{red}J)$, alors $\underline{\operatorname{Hom}}(J,M)$ est $\pi$-accessible, plus précisément on a une équivalence de catégories

\[
\underline{\operatorname{Hom}}(J,M)\xleftarrow{\approx}\operatorname{Ind}_{\pi}(\underline{\operatorname{Hom}}(J,M_{\pi}))\tag*{(4.9.5.1)}
\]
% archive sheet number: 303
